\section{Current Source-Owned Soft Contact Algorithm}

This section describes the current Python dynamics algorithm as implemented in
\texttt{Base.py} and \texttt{NeoDynamics.py}.  The algorithm is intended to be
the reference behavior for the future GLSL implementation.

The important modeling goals are:
\begin{itemize}
    \item particles are soft and may penetrate during contact,
    \item contacts are purely repulsive,
    \item tangential motion should pass through contact without artificial
          friction,
    \item wall contacts should use the same geometry path as particle-particle
          contacts by using ghost particles,
    \item each source particle may write only its own particle state and its
          own source-owned contact state.
\end{itemize}

\subsection{Step Order}

For each simulation substep, \texttt{Base.step} performs the following sequence:
\begin{enumerate}
    \item Begin the integration step and prepare the active position buffer.
    \item Detect particle-particle and particle-wall contacts.
    \item Update source-owned particle contact state.
    \item Update source-owned wall ghost contact state.
    \item Predict and apply source-owned collision velocities.
    \item Move particles using the updated velocities.
    \item Detect contacts again at the moved positions.
    \item Update wall contact state again.
    \item Clip rebound overlaps that exceed the stored zero-velocity geometry.
    \item End the integration step and update diagnostic internal momentum.
\end{enumerate}

The second contact detection after movement is important because overlap
geometry may change during the substep.  The clipping pass is intentionally
geometric and source-owned.

\subsection{Contact Geometry}

For source particle \(i\) and target particle \(j\), the center distance is
\[
    d_{ij} = \lVert \mathbf{x}_j - \mathbf{x}_i \rVert .
\]
A contact exists when
\[
    d_{ij} < r_i + r_j .
\]
The source contact normal points from source to target:
\[
    \mathbf{n}_{ij}
    =
    \frac{\mathbf{x}_j - \mathbf{x}_i}{d_{ij}} .
\]
If the centers coincide, the implementation chooses a stable fallback normal.

The relative normal velocity is
\[
    u_{ij}
    =
    (\mathbf{v}_j - \mathbf{v}_i) \cdot \mathbf{n}_{ij}.
\]
For particle-particle contacts, \(u_{ij}<0\) means compression and
\(u_{ij}>0\) means separation.  The contact phase stored in source-owned
contact state is either \texttt{compression} or \texttt{rebound}.

\subsection{Source-Owned Contact State}

Each particle owns an array \texttt{gcs} of contact-state records.  A record is
owned by the source particle and stores either a real particle target index or a
wall target of the form \texttt{("wall", wall\_flag)}.

A contact-state record stores:
\begin{itemize}
    \item the contact type, target id, and active flag,
    \item the phase, \texttt{compression} or \texttt{rebound},
    \item the motion mode, \texttt{accumulating}, \texttt{releasing}, or
          \texttt{blocked},
    \item the first-contact normal,
    \item first-contact source and target velocities,
    \item the zero-velocity overlap area \(A_0\),
    \item the zero-velocity center distance \(d_0\),
    \item diagnostic internal momentum values.
\end{itemize}

The source-owned rule is strict: a source particle may read target particle
state, but it must not write target position, target velocity, or target-owned
contact state.

\subsection{Wall Ghost Contacts}

Walls are represented by virtual ghost particles.  For a real particle of
radius \(r\), the ghost has the same radius and copies the real particle's
contact material parameters.  The ghost is fixed, has zero velocity, and is
given a very large effective mass.

The ghost perimeter is moved slightly inside the visible wall by
\[
    s = \frac{1}{8}r .
\]
For the right wall, for example, with wall coordinate \(x_{\max}\), the ghost
center is
\[
    \mathbf{x}_g = (x_{\max} + r - s,\; y_i).
\]
The other walls follow the same pattern.  The ghost is outside the domain, but
its perimeter overlaps the wall line by the skin distance \(s\).  Contact is
then evaluated as an ordinary equal-radius particle-particle overlap between
the real particle and the ghost.

This makes the wall contact law geometrically similar to particle-particle
contact.  The wall is still fixed because the ghost is not integrated and any
ghost velocity calculated by the pair model is discarded.

\subsection{Zero-Velocity Geometry}

The Neo force-law model computes a zero-velocity state for each contact.  The
zero-velocity point is the point of maximum penetration: compression ends and
rebound begins.

The stored zero geometry is:
\[
    A_0 = \text{zero-velocity overlap area},
\]
\[
    d_0 = \text{zero-velocity center distance}.
\]
The model uses overlap area as the progress coordinate.  During compression,
overlap increases toward \(A_0\).  During rebound, overlap decreases away from
\(A_0\).

For contacts that exist at frame zero, the current overlap is treated as being
on the rebound side of the collision history.  The model reconstructs a
zero-velocity state from the starting overlap and the incoming contact
momentum.

\subsection{Velocity Prediction}

Velocity prediction is source-owned.  Each source particle predicts only its
own next velocity.  For each active particle contact and wall contact, the
source particle asks \texttt{NeoDynamics} for a compression or rebound velocity
at the current overlap.

For a particle-particle contact, the prediction path is:
\begin{enumerate}
    \item read the source particle, target particle, contact normal, overlap
          area, and center distance,
    \item update the source-owned phase and motion mode,
    \item construct a temporary two-particle Neo model,
    \item compute the current-overlap compression or rebound velocity report,
    \item apply only the source particle's predicted normal component,
    \item preserve the source particle's current tangential component.
\end{enumerate}

Only the normal component is replaced.  Let
\[
    v_n = \mathbf{v}_i \cdot \mathbf{n}_{ij},
\]
and let \(v_n^\ast\) be the predicted normal component from the Neo report.
The source velocity update is
\[
    \mathbf{v}_i^\ast
    =
    \mathbf{v}_i
    +
    (v_n^\ast - v_n)\mathbf{n}_{ij}.
\]
The tangential component is not modified by the contact response.

When a source particle has multiple active contacts, the code applies predicted
normal components sequentially to the same source velocity accumulator.  This
is source-owned and GPU-friendly, but it is not a full coupled cluster solve.
This is the current place where difficult glancing multi-contact behavior can
appear.

\subsection{Rebound Clamp}

The current Python path still includes the parameter
\texttt{geo\_rebound\_min\_fraction}.  Its default value is \(0.02\).  During
rebound, the Neo model may clamp the predicted rebound velocity so that the
outward normal speed is not below a small fraction of the original closing
speed.

This clamp is a numerical guard, not a fundamental physical rule.  Larger
values can force contacts to separate quickly, but they also make soft contacts
look too hard.  The intended physical model should not rely on this parameter
as the primary reason particles separate.

\subsection{Geometric Overlap Clipping}

After velocity prediction and movement, rebound contacts are geometrically
clipped if they penetrate deeper than their stored zero-velocity geometry.  For
particle-particle contacts, if
\[
    A > A_0,
\]
the source particle is moved so that the pair center distance returns toward
\[
    d_0 .
\]
The correction is source-owned and mass-weighted.  If
\[
    c = d_{ij} - d_0,
\]
then the source particle moves by
\[
    \Delta \mathbf{x}_i
    =
    c
    \frac{m_j}{m_i + m_j}
    \mathbf{n}_{ij}.
\]
The target particle is not written by this source invocation.  When both source
particles run their own source-owned correction, the combined effect is the
usual mass-weighted pair correction without either invocation writing its
target.

For wall contacts, the real particle is moved along the wall normal so that the
real-particle to ghost-particle center distance matches the stored
zero-velocity distance.  The wall ghost is never moved or integrated.

\subsection{Internal Momentum Accounting}

The model reports internal momentum as a diagnostic quantity.  Internal normal
momentum is computed from first-contact normal momentum and current normal
momentum.  It is accumulated into per-particle diagnostic fields:
\[
    I_x = I\,n_x,
    \qquad
    I_y = I\,n_y .
\]
These fields are useful for understanding how much momentum is stored in active
contacts, but they are not independent forces applied to particles.

The UI reports:
\begin{itemize}
    \item starting particle momentum,
    \item current particle momentum,
    \item internal contact momentum,
    \item returned wall momentum diagnostic,
    \item corrected momentum difference.
\end{itemize}

\subsection{Known Failure Mode}

The current soft source-owned algorithm can still produce a glancing-contact
artifact.  In \texttt{TwoParticleStuckWall.cfg}, particle 0 can remain in wall
contact while particle 1 remains in particle-particle contact with particle 0.
The wall response changes particle 0's motion, while the particle-particle
contact remains soft and overlapped.  As the particles slide, the
particle-particle normal rotates.  Visually this can look like particle 1 is
rotating around particle 0 or like an artificial attraction.

This is not caused by an explicit attractive force.  It is caused by a
long-lived soft overlap whose normal changes while velocity prediction is still
being applied contact-by-contact.

The attempted hard geometric response eliminated the rotation but also removed
soft penetration from simple cases such as \texttt{TwoParticleHorizontal.cfg}.
Therefore the velocity law must remain soft, while the geometric correction and
future multi-contact handling must prevent nonphysical orbit-like behavior.

\subsection{Implications For GLSL}

The GLSL implementation should preserve the source-owned rule:
\begin{itemize}
    \item one invocation owns one source particle,
    \item it may read neighboring particles and wall ghosts,
    \item it may write only the source particle and source-owned contact state,
    \item pair-level data is diagnostic, not an independent owner of response.
\end{itemize}

The GLSL code should not implement pair-owned target writes.  If a future
multi-contact solver is needed, it must still be expressed as source-owned
updates or as a separate staged process with explicit ownership rules.
